\section*{Программа экзамена.}
\begin{enumerate}
    \item Слабая-* топология в сопряжённом пространстве топологического векторного пространства. Теорема об общем виде слабо-* непрерывного линейного функционала над сопряжённым пространством.
    
    \Cref{seq:weak-star-topology} \nameref{seq:weak-star-topology}. 

    \Cref{lemma:about-continious-funcitonals} \nameref{lemma:about-continious-funcitonals}.

    \item Теорема Банаха-Алаоглу о слабой-* компактности поляры окрестности нуля топологического векторного пространства.
    
    \Cref{seq:banach-alaogly}. \nameref{seq:banach-alaogly}.

    \item Метризуемость слабой-* топологии на шаре в сопряжённом пространстве сепарабельного нормированного пространства и секвенциально слабая-* компактность замкнутого шара в сопряжённом пространстве.
    
    \Cref{seq:metrizability}. \nameref{seq:metrizability}. Теорема 2.3.1, утверждение 2.

    \item Отсутствие метризуемости слабой-* топологии в сопряжённом пространстве бесконечномерного банахова пространства.
    
    \Cref{seq:unmetrizability}. \nameref{seq:unmetrizability}. Теорема 2.2.1, утверждение 2.

    \item Теорема об отделимости двух непересекающихся выпуклых подмножеств топологического векторного пространства.
    
    \Cref{def:minkovsky}. \Cref{prop:minkovsky_good}. \Cref{th:separability_first} <<\nameref{th:separability_first}>>.

    \item Теорема об отделимости точки от несодержащего её выпуклого замкнутого подмножества локально-выпуклого топологического векторного пространства.
    
    \Cref{th:separability_second} <<\nameref{th:separability_second}>>.

    \item Слабая топология в нормированном пространстве, её сравнение с нормированной топологией. Теорема Мазура о слабой замкнутости выпуклого сильно замкнутого подмножества нормированного пространства.
    
    \Cref{seq:weak-topology} \nameref{seq:weak-topology}. \Cref{prop:infdim-subspace} и следствие из него. \Cref{th:mazur}. Следствия из теоремы Мазура.

    \item Отсутствие метризуемости слабой топологии в бесконечномерном нормированном пространстве.
    
    \Cref{seq:unmetrizability}. \nameref{seq:unmetrizability}. Теорема 2.2.1, утверждение 1.

    \item Метризуемость слабой топологии на шаре нормированного пространства, сопряжённое пространство которого сепарабельно.
    
    \Cref{seq:metrizability}. \nameref{seq:metrizability}. Теорема 2.3.1, утверждение 1.

    \item Сепарабельность нормированного пространства, сопряжённое которого сепарабельно.
    
    \Cref{prop:conjugate_separability}.

    \item Теорема Эберлейна-Шмульяна о слабой секвенциальной компактности слабого компакта в нормированном пространстве.
    
    \Cref{seq:eberlein-smulyan} <<\nameref{seq:eberlein-smulyan}>>.

    \item Слабая компактность замкнутого единичного шара в рефлексивном нормированном пространстве. Существование метрической проекции точки на выпуклое замкнутое подмножество рефлексивного пространства.
    
    \Cref{seq:banach-alaogly}. Следствие из теоремы Банаха-Алаоглу.
    \ref{cor:metric-projection}. Является следствием из утверждения 14 пункта программы.

    \item Слабая компактность замкнутого единичного шара в нормированном пространстве как критерий его рефлексивности.
    
    \Cref{cor:reflectia-criterion}.

    \item Теорема о существовании минимума на выпуклом замкнутом ограниченном подмножестве рефлексивного пространства сильно непрерывного выпуклого вещественного функционала.
    
    \Cref{cor:separability-corollaries} и дальше ещё следствия три.

    \item Теорема Рисса-Фреше о представлении сопряжённого гильбертового пространства. Рефлексивность гильбертового пространства.
    
    \Cref{seq:riesz-frechet} <<\nameref{seq:riesz-frechet}>>.

    \Cref{th:hilbert-reflexivity}. Note: доказательство самопальное, в лекциях есть каноническое (конец 5 лекции).

    \item Аннуляторы $M^\perp$ и ${}^\perp N$ подпространств $M \subset X$ и $N \subset X^*$ для нормированного пространства $X$. Вычисление $^\perp(M^\perp)$ и $(^\perp N)^\perp$. Слабая-* плотность в $X^{**}$ подпространства $F(X)$, где $F\colon X \to X^{**}$ -- каноническая изометрия Банаха.
    

    \Cref{def:annulators}.
    Вычисление двойных аннуляторов: \Cref{cor:annulator_equations}.
    \Cref{th:goldstein}, теорема Голдстайна.

    \item Равносильность рефлексивности банахова пространства и рефлексивности его сопряжённого.

    \Cref{th:refref}.  Note: доказательство самопальное.

    \item Оператор, сопряжённый к $A \in \LL(X, Y)$. Равенства $^\perp (\ker A^*)$ сильному замыканию $\im A$ и $(\ker A)^\perp$ слабому-* замыканию $\im A^*$. Равенство $(\ker A)^\perp = \im A^*$ при условии замкнутости $\im A$ и полноты пространств $X$ и $Y$.
    
    \Cref{seq:conjugate-operators} \nameref{seq:conjugate-operators}. \Cref{th:fredhgolm}. \Cref{th:closability}.

    \item Теорема о замкнутости $\im A$ для $A \in \LL(X, Y)$ при условиях замкнутости $\im A^*$ и полноты пространства $X$.
    
    \Cref{th:another-closability}. После теоремы идёт примечание о существенности банаховости.

    \item Конечномерность $\ker A_\lambda$ и замкнутость $\im A_\lambda$ для компактного оператора $A \in \LL(X)$ и нетривиального $\lambda \in \Cm$, где $X$ -- банахово.
    
    \Cref{th:finite-dimension-kernel}.

    \item Эквивалентность компактности $A \in \LL(X)$ и $A^* \in \LL(X^*)$, где $X$ -- банахово пространство.
    
    \Cref{prop:shauder-theorem}.

    \item Равенство размерностей $\ker A_\lambda$ и $\ker (A_\lambda)^*$ для компактного оператора $A \in \LL(H)$ и $\lambda \in \Cm \setminus \{0\}$, где $H$ -- гильбертово пространство.
    
    \Cref{th:dim-ker-equality}. И ещё предыдущая теорема, полагаю.

    \item Критерий Фредгольма разрешимости уравнения $A_\lambda x = y$ для компактного оператора $A \in \LL(X)$ и нетривиального $\lambda \in \Cm$, где $X$ -- банахово пространство. Альтернатива Фредгольма.
    
    Альтернатива Фредгольма идёт после теоремы о равенстве размерностей ядер.
    \begin{notet}[Критерий Фредгольма]
        \[
        A_\lambda x=y \text{ разрешимо}
        \Longleftrightarrow
        y\in \operatorname{Im}A_\lambda
        \Longleftrightarrow
        y\in {}^\perp(\ker A_\lambda)
        \]
        то есть
        \[
        A_\lambda x=y \text{ разрешимо}
        \Longleftrightarrow
        f(y)=0\quad \forall f\in\ker A_\lambda,
        \]
        где $A$ -- компактный оператор и $\lambda \in \Cm \setminus \{0\}$. Доказательство -- применение аннуляторных соотношений и замкнутости $\im A_\lambda$ для компактного оператора при нетривиальном $\lambda$.
    \end{notet}

    \item Теорема о собственных числах компактного оператора и теорема о спектре компактного оператора в банаховом пространстве.
    
    \Cref{seq:compact_spectrum}. Теорема о спектре компактного оператора, собственно.

    \item Теорема о вещественности спектра самосопряжённого оператора в гильбертовом пространстве. Пустота остаточного спектра самосопряжённого оператора.
    
    \Cref{seq:selfadjoint}. Утверждения 1 -- 2 теоремы о спектре самосопряжённого оператора.
    Пустота остаточного спектра идёт в формате заметки после доказательства теоремы.

    \item Критерий вложения вещественного числа в спектр самосопряжённого оператора в гильбертовом пространстве. Локализация спектра самосопряжённого оператора.
    
    \Cref{seq:selfadjoint}. Утверждения 3 -- 5 теоремы о спектре самосопряжённого оператора.

    \item Теорема Гильберта-Шмидта для компактного самосопряжённого оператора в гильбертовом пространстве.
    
    \Cref{seq:gilbert-shmidt}. \nameref{seq:gilbert-shmidt}

    \item Открытость резольвентного множества, непустота и компактность спектра непрерывного оператора в банаховом пространстве.
    
    \Cref{seq:spectrum}. \nameref{seq:spectrum}


    \item Теорема о спектральном радиусе непрерывного оператора в банаховом пространстве. Критерий равенства спектрального радиуса норме непрерывного оператора.
    
    \Cref{seq:gelfand}. \nameref{seq:gelfand}.

\end{enumerate}